\documentclass[12pt,a4paper]{article}
\usepackage[utf8]{inputenc}
\usepackage{mathptmx}
\usepackage[margin=1in]{geometry}
\usepackage{amsmath,amssymb}
\usepackage{graphicx}
\usepackage{float}
\usepackage{setspace}
\usepackage[round]{natbib}
\usepackage[colorlinks=true,linkcolor=blue,citecolor=blue,urlcolor=blue]{hyperref}
\onehalfspacing

\title{\textbf{The Gravitational Shadow: Random Matrix Theory Reveals Universal Level Repulsion in Planetary Orbital Architecture}\\[0.4em]
\large A Local-Spacetime Repulsion Law from the Solar System to 320 Kepler Systems}
\author{Ruqing Chen\\ GUT Geoservice Inc., Montreal, Canada\\ \texttt{ruqing@hotmail.com}}
\date{June 2026}

\begin{document}
\maketitle

\begin{abstract}
\noindent A random matrix theory (RMT) program has shown that single-source, charge-and-release systems across the Earth sciences exhibit Wigner--Dyson level repulsion in the timing of their events, while superpositions of many independent sources collapse to Poisson. Here we carry the program into celestial mechanics, where the relevant ``charge-and-release'' process is gravitational accretion: a forming planet clears its feeding zone (its Hill sphere and resonance-swept neighborhood), casting a \emph{gravitational shadow} that forbids other bodies from occupying nearby orbits. We test whether this clearing imprints level repulsion on the logarithmic spacing of orbital semi-major axes, $\Delta\ln a$. Two methodological traps are handled explicitly: the Titius--Bode trap (planetary $\ln a$ is near-geometric, so a naive spacing ratio is inflated by periodicity, not repulsion) is removed by local unfolding of each system; and the resonance trap (mean-motion-resonant chains lock into exactly geometric spacing) is flagged by a post-unfolding coefficient of variation below 0.15. After unfolding, the Solar System's eight planets show genuine repulsion ($\langle r\rangle=0.632$), while the Galilean moons are correctly flagged as resonance-locked (CV~$=0.09$, Laplace 1:2:4) and excluded. Pooling 320 Kepler/transit multi-planet systems (1130 planets), each unfolded in its own frame, yields $\langle r\rangle=0.674$ (95\% CI [0.68, 0.72]), GUE-class repulsion that is \emph{stronger} than the Solar System and stronger than any Earth system in the program. This contradicts the naive prediction of Poisson collapse: pooling does not randomize, because each planetary system is \emph{internally} repulsive and we preserve that structure by unfolding each system in its own frame. The collapse is recovered only when the per-system frames are destroyed: globally mixing all 1130 planets into a single catalog drives $\langle r\rangle$ down to 0.437, and an independent-clock null gives 0.396. Level repulsion in orbital architecture is therefore a \emph{local-spacetime} property: within a single gravitational well it is universal and GUE-strong, but mixing many independent wells---different stellar masses, formation histories, and dynamical clocks---erases it. The above-GUE rigidity characterizes the dynamically cold, tightly packed systems of the Kepler dichotomy, plausibly the maximum-packing limit of planetary formation. The gravitational shadow is the spatial analog of the temporal charge-and-release mechanism, and $\langle r\rangle$ measures the rigidity of orbital clearing.

\vspace{0.8em}
\noindent\textbf{Keywords:} random matrix theory; level repulsion; orbital architecture; exoplanets; Kepler; gravitational clearing; Hill sphere; dynamical packing; Titius--Bode; spectral superposition

\vspace{0.4em}
\noindent\textbf{Code and data:} \url{https://github.com/Ruqing1963/orbital-rmt-gravitational-shadow}

\vspace{0.2em}
\noindent\textbf{Companion studies:} stratigraphy \url{https://zenodo.org/records/20774581}; seismotectonics \url{https://zenodo.org/records/20768130}; mantle plumes \url{https://zenodo.org/records/20768420}; metallogeny \url{https://zenodo.org/records/20768849}; evolution \url{https://zenodo.org/records/20783763}; hydrogeology \url{https://zenodo.org/records/20780389}; solar flares \url{https://zenodo.org/records/20784967}
\end{abstract}

\section{Introduction}

A random matrix theory (RMT) program has tested a single hypothesis across many physical systems: that the spacing of events generated by a single, long-memory ``charge-and-release'' source exhibits level repulsion of the Wigner--Dyson form, while the superposition of many independent sources collapses to Poisson \citep{chen2026plume,chen2026metal,chen2026strat,chen2026evo,chen2026hydro,chen2026solar}. The evidence has been consistent across the Earth and stellar sciences: mantle plumes ($\langle r\rangle=0.630$), ore-forming systems (0.57--0.71), cyclostratigraphic spacings, mass extinctions and radiations (0.62), a geyser (0.738), glacial outburst floods (0.550), and---at the low-depletion limit---solar active regions (Poisson). In every case the diagnostic has been temporal: the waiting times \emph{between} events.

Here we take the program into a new domain, where the relevant variable is not temporal but \emph{spatial}, and the mechanism is gravity. When a planet forms, it accretes and scatters the planetesimals in its feeding zone, clearing an annulus of width set by its Hill sphere and by mean-motion resonances. This cleared annulus is a \emph{gravitational shadow}: a region of orbital space from which other surviving bodies are excluded. The clearing is the spatial counterpart of charge-and-release---the disk's local solid mass is the reservoir, and accretion is the release that empties it. If the Earth-systems hypothesis generalizes, the cleared shadow should forbid bodies from occupying neighboring orbits, producing level repulsion in the \emph{spacing} of orbital semi-major axes, just as a recharge time forbids closely spaced events in time.

We test this on the logarithmic spacing $\Delta\ln a$ of orbital semi-major axes, the natural variable because planetary systems are approximately scale-invariant (period ratios, not absolute separations, govern stability). We study two targets. Target A is the single-source regime: the Solar System's eight planets and the Galilean moon system, each a single gravitational well. Target B is the would-be multi-source regime: 320 Kepler and transit/radial-velocity multi-planet systems pooled together. The original expectation, by analogy with the Earth-systems program, was that Target B would collapse to Poisson. We show that it does not---and that the reason it does not is itself the deepest result of this study.

\section{The Gravitational Shadow Mechanism}

\subsection{Clearing as spatial charge-and-release}

In the temporal program, repulsion requires that each release substantially depletes a single source, forcing a recharge interval before the next comparable event and thereby forbidding short waiting times. The spatial analog is immediate. A forming planet depletes the solid material in its feeding annulus; the cleared zone forbids the long-term survival of another body at a nearby semi-major axis, because such a body would lie within the first planet's Hill sphere or a destabilizing resonance and be ejected or accreted on dynamical timescales. The cleared annulus is a hard minimum spacing, the spatial equivalent of a recharge time. Orbital architecture should therefore show level repulsion in $\Delta\ln a$: planets repel in orbital space exactly as charge-and-release events repel in time.

\subsection{Why $\ln a$ is the canonical dynamical variable}

The choice of $\ln a$, rather than the absolute separation $\Delta a$, is dictated by the dynamics, not by convenience. In celestial mechanics, both long-term Hill stability and the locations of mean-motion resonances are governed by \emph{period ratios} $P_{i+1}/P_i$, not by absolute spacings: two planets are dynamically close when their periods are commensurable or their period ratio approaches the overlap of first-order resonances, regardless of how many AU separate them. By Kepler's Third Law, $P^2\propto a^3$, so $\ln P = \tfrac{3}{2}\ln a + C$, and therefore
\begin{equation}
\ln\!\left(\frac{P_{i+1}}{P_i}\right) = \frac{3}{2}\,\bigl(\ln a_{i+1}-\ln a_i\bigr) = \frac{3}{2}\,\Delta\ln a .
\end{equation}
Equal logarithmic spacing in semi-major axis corresponds \emph{exactly} to equal logarithmic spacing in orbital period, and the period-ratio spectrum that controls stability is, up to the constant factor $3/2$, the $\Delta\ln a$ spectrum we analyze. Thus $\ln a$ is the canonical variable for expressing the gravitational shadow: it is the coordinate in which the dynamical clearing acts uniformly. Using $\Delta a$ instead would mix the strong radial gradient in absolute spacing (planets are farther apart farther out) into the statistics, manufacturing a spurious trend that has nothing to do with repulsion.

\subsection{The effective local frame}

As in the temporal program, the repulsion is a property of an \emph{effective local frame}, not of any global ledger. Each planetary system has its own characteristic scale: its stellar mass sets the absolute $a$ values, its disk mass and surface-density profile set the mean planet spacing, and its dynamical age sets how completely the architecture has relaxed. Repulsion lives in the dimensionless, locally unfolded spacing within a single system. This is the spatial counterpart of the asperity in the seismotectonic case \citep{chen2026fault}: the relevant reservoir is the one local, isolated, dynamically lockable parcel, not the entire disk or galaxy. We will see that this locality is not a technical nicety but the physical crux: it is precisely what distinguishes the universal within-system repulsion from the Poisson collapse that follows when independent local frames are mixed.

\section{Data and Methods}

\subsection{Orbital data}

For Target A we use textbook semi-major axes: the eight Solar System planets ($a=0.387$--30.07~AU) and the four Galilean moons of Jupiter. For Target B we query the NASA Exoplanet Archive Planetary Systems Composite Parameters table (\texttt{pscomppars}) for all systems with three or more planets having measured semi-major axes:
\begin{quote}\footnotesize
\texttt{select hostname, pl\_name, pl\_orbsmax, sy\_pnum from pscomppars where pl\_orbsmax>0 and sy\_pnum>2}
\end{quote}
This returns 1130 planets in 341 host systems, of which 320 systems have three or more distinct, valid semi-major axes (117 with four or more, 38 with five or more). The sample is overwhelmingly dominated by Kepler and other transit/radial-velocity discoveries in the solar neighborhood; all systems are Galactic (no extragalactic planet has a resolved orbital catalog).

\subsection{The Titius--Bode trap and local unfolding}

Planetary semi-major axes follow an approximately geometric progression ($a_{n}\approx a_0 k^n$, the Titius--Bode regularity), so $\ln a$ is approximately \emph{equally spaced}. An equally spaced sequence has spacing ratio $\langle r\rangle\to1$ purely from periodicity, masquerading as super-rigid repulsion. The raw Solar System gives $\langle r\rangle=0.606$ and the raw Galilean moons $\langle r\rangle=0.911$, both contaminated by this geometric trend. We remove it by \emph{local unfolding}: for each system we fit the cumulative count $N(\ln a)$ with a smooth function (a linear fit for systems with $\le5$ bodies, a smoothing spline for larger systems), evaluate the local mean density, and rescale the spacings to unit mean. This is the same unfolding used to remove the Sepkoski-curve trend in the evolution study \citep{chen2026evo}. After unfolding, a pure geometric sequence yields $\mathrm{CV}\to0$ (no repulsion, only periodicity), genuine GOE repulsion yields $\mathrm{CV}\approx0.52$, and Poisson yields $\mathrm{CV}\approx1.0$.

\subsection{The resonance trap}

Mean-motion resonant chains (for example the Galilean Laplace 1:2:4 resonance) lock bodies into exactly geometric spacing. After unfolding, such a system collapses to $\mathrm{CV}\to0$: its high $\langle r\rangle$ reflects resonant periodicity, not Wigner repulsion. We flag any system with post-unfolding $\mathrm{CV}<0.15$ as resonance-locked and decline to interpret its $\langle r\rangle$ as RMT repulsion. This flag fires correctly on the Galilean system and on no Solar System planetary case.

\subsection{RMT diagnostics}

For each unfolded spacing set we compute the spacing ratio $\langle r\rangle$, the coefficient of variation CV, the Brody parameter $\hat\beta$, Kolmogorov--Smirnov distances to Poisson, GOE, and GUE, and 95\% bootstrap confidence intervals on $\langle r\rangle$. Reference values are Poisson $\langle r\rangle=0.386$, GOE $0.531$, GUE $0.603$.

\section{Results}

\subsection{Target A: single gravitational wells}

The Solar System's eight planets, after local unfolding, give $\langle r\rangle=0.632$ with $\mathrm{CV}=0.344$ (Figure~\ref{fig:repulsion}a). The resonance flag does not fire (CV well above 0.15), so this is genuine level repulsion: the residual orbital spacings, with the Titius--Bode geometric trend removed, are more uniform than random, in the GOE--GUE range. Gravitational clearing has imprinted repulsion on the spacing of the planets.

The Galilean moons give a raw $\langle r\rangle=0.911$, but after unfolding the coefficient of variation is only 0.094: the resonance flag fires (Table~\ref{tab:main}). This near-zero CV is the signature of the Laplace 1:2:4 mean-motion resonance, which locks Io, Europa, and Ganymede into an essentially exact geometric chain. We therefore report the Galilean system as resonance-locked and \emph{not} as evidence of RMT repulsion: its rigidity is periodic, not Wigner. This is a deliberate negative call---the apparent ``super-GUE'' value is an artifact, and reporting it as repulsion would be a methodological error of the kind the unfolding and resonance flags are designed to catch.

Target A's single-system samples are small (the Solar System provides only six spacing ratios), so they are qualitative; the statistical weight comes from Target B.

\subsection{Target B: pooling 320 systems does not collapse to Poisson}

Pooling the locally unfolded spacings of all 320 multi-planet systems (810 spacings) gives $\langle r\rangle=0.674\pm0.008$ with a 95\% bootstrap CI of [0.68, 0.72], $\mathrm{CV}=0.341$, and a Brody parameter $\hat\beta=2.11$ (Figure~\ref{fig:repulsion}b). The KS test prefers GUE ($D=0.079$) over GOE ($D=0.134$) and decisively over Poisson ($D=0.340$). This is GUE-class repulsion---\emph{stronger} than the Solar System single system, and stronger than any Earth system in the program.

This flatly contradicts the naive prediction. By analogy with the Earth-systems program, pooling hundreds of independent gravitational wells was expected to collapse to Poisson via spectral superposition. It does the opposite. The reason is structural and is the central physical result of this study (Section~5.1): each planetary system is \emph{itself} internally repulsive, and because we unfold each system in its own frame before pooling, we preserve that internal repulsion rather than averaging it away. Pooling many internally repulsive systems, each normalized to its own frame, does not randomize; it reveals the \emph{universality} of the internal repulsion across 320 independent stars.

\begin{table}[H]
\centering
\caption{RMT diagnostics for orbital architecture. The Solar System and the pooled Kepler ensemble both show GUE-class repulsion in unfolded orbital spacing. The Galilean system is flagged resonance-locked and excluded. The collapse controls (Section~4.6) show that destroying the per-system frames recovers Poisson.}
\label{tab:main}
\small
\begin{tabular}{lccccl}
\toprule
Target & $n$ & $\langle r\rangle$ & 95\% CI & CV & Class \\
\midrule
Solar System (8 planets, unfolded) & 7 & 0.632 & [0.55, 0.87] & 0.34 & GOE--GUE \\
Galilean moons (4, unfolded) & 3 & (0.911) & --- & 0.09 & \emph{resonance-locked} \\
\textbf{Kepler 320 systems (aligned pool)} & \textbf{810} & \textbf{0.674} & \textbf{[0.68, 0.72]} & 0.34 & \textbf{GUE} \\
\midrule
Kepler global mix (frames scrambled) & 827 & 0.437 & [0.41, 0.45] & 0.86 & Poisson \\
Poisson null (independent clocks) & 810 & 0.396 & [0.36, 0.41] & 1.05 & Poisson \\
\bottomrule
\end{tabular}
\end{table}

\begin{figure}[H]
\centering
\includegraphics[width=\textwidth]{figs/fig1_orbital_repulsion.pdf}
\caption{Orbital level repulsion. (a) The Solar System's eight planets after Titius--Bode unfolding: the spacing distribution sits in the GOE--GUE range ($\langle r\rangle=0.632$), genuine repulsion. (b) The 320 pooled Kepler/transit multi-planet systems ($n=810$ unfolded spacings): the data (blue KDE) track GUE (purple) rather than the predicted Poisson (dotted), with $\langle r\rangle=0.674$. Orbital architecture is GUE-repulsive at both the single-system and ensemble levels.}
\label{fig:repulsion}
\end{figure}

\subsection{The collapse test: repulsion is a local-spacetime property}

If the pooled GUE signal reflects internal, per-system repulsion, then destroying the per-system frames should recover the Poisson collapse. We test this directly with two controls (Figure~\ref{fig:collapse}). First, we \emph{globally mix}: all 1130 planets from all systems are thrown into a single catalog and unfolded \emph{once}, as if they shared one cosmic ruler. This scrambles the independent stellar masses, disk properties, and dynamical clocks of the contributing systems. The result is $\langle r\rangle=0.437$ ($\mathrm{CV}=0.86$), collapsed far below the aligned value and toward Poisson. Second, we construct an \emph{independent-clock null} by replacing each system with the same number of planets placed at random in $\ln a$ over its span, then pooling; this gives $\langle r\rangle=0.396$ ($\mathrm{CV}=1.05$), Poisson as expected for genuinely independent random sources.

The contrast is decisive. The same 320 systems give GUE repulsion when each is read in its own frame (0.674) but collapse toward Poisson when their frames are mixed into one (0.437). Level repulsion in orbital architecture is therefore a \emph{local-spacetime} property: it exists within each single gravitational well, where one stellar mass and one disk set a common scale, and it is destroyed when many independent wells---each with its own scale, formation history, and dynamical clock---are superposed without alignment. This is the spatial, gravitational realization of the program's central dichotomy: single source repels, mixed independent sources collapse to Poisson. The novelty is that here ``single source'' means ``single spacetime frame,'' and the collapse is recovered not by adding more systems but by erasing the frames that align them.

\begin{figure}[H]
\centering
\includegraphics[width=\textwidth]{figs/fig2_collapse_test.pdf}
\caption{The collapse test. (a) The aligned pool (each system unfolded in its own frame) is GUE-repulsive ($\langle r\rangle=0.674$). (b) Globally mixing all 1130 planets into one catalog scrambles the independent per-system frames and collapses the distribution toward Poisson ($\langle r\rangle=0.437$). (c) An independent-clock null (each system replaced by random $\ln a$ points) gives Poisson ($\langle r\rangle=0.396$). Repulsion survives within a shared local frame and is destroyed by mixing independent ones.}
\label{fig:collapse}
\end{figure}

\section{Discussion}

\subsection{Why pooling does not collapse: a correction to the naive superposition picture}

The Earth-systems program established that superposing many independent sources collapses to Poisson. It is tempting to read this as ``many systems $\Rightarrow$ Poisson,'' and to expect 320 planetary systems to wash out to randomness. Our result corrects this. Spectral superposition collapses to Poisson only when the superposed sources are themselves \emph{independent random (Poisson) clocks with no shared scale}. The Earth-systems superpositions satisfied this---mixing earthquakes from unrelated faults, or eruptions from unrelated volcanoes, combines sources with no common ruler and random relative phase. Planetary systems do not satisfy it. Each system is internally repulsive (dynamical packing, orbital clearing, long-term stability all enforce minimum spacings), so it is not a random clock; and when we unfold each system in its own frame, we explicitly align them to a common dimensionless scale before pooling. Pooling aligned, internally repulsive systems preserves and reveals their shared repulsion. The collapse returns only when we remove the alignment (global mix) or remove the internal repulsion (Poisson null). The lesson is that the controlling variable is not the number of sources but whether they share a single local frame.

\subsection{$\langle r\rangle$ as a probe of clearing rigidity}

As in the temporal program, $\langle r\rangle$ becomes a quantitative instrument. A planetary architecture with $\langle r\rangle$ near GUE has been dynamically sculpted to near-maximal rigidity: its planets occupy the most uniformly repelling configuration consistent with stability, the spatial analog of a fully depleted, hard-recharging reservoir. The pooled Kepler value (0.674) slightly exceeds GUE, indicating that surviving multi-planet systems are, if anything, \emph{over}-packed into rigid configurations---a natural consequence of the stability filter, which preferentially destroys the systems whose spacings were too irregular (too close to Poisson) to survive. The diagnostic thus reads out the strength of gravitational clearing and the severity of the long-term stability filter that has acted on each surviving system.

\subsection{The packing limit and the Kepler dichotomy}

The above-GUE rigidity connects naturally to the \emph{Kepler dichotomy}, the observation that planetary systems divide into dynamically ``cold'' (compact, coplanar, multi-transiting) and dynamically ``hot'' (sparse or single, more inclined and eccentric) populations. Our 320-system sample, by construction (three or more detected planets with measured semi-major axes), is dominated by the cold, tightly packed population. Two effects act in the same direction here, and intellectual honesty requires naming both. First, a genuine \emph{detection bias}: transit and radial-velocity surveys preferentially recover compact, coplanar systems, so the sample over-represents dynamically cold architectures. Second, a genuine \emph{physical filter}: among the systems that do form compactly, only those sculpted to the most uniform, maximally repelling spacings survive Gyr of dynamical evolution; the rest go unstable and are pruned. The observed $\langle r\rangle=0.674$ therefore characterizes the surviving, dynamically cold systems of the Kepler dichotomy, which represent something close to the \emph{maximum-packing limit} of planetary formation: architectures sculpted to the brink of instability, forced into the most uniform repulsive configuration available. We propose that GUE-plus rigidity is the statistical signature of this packing limit---a kind of dynamical crystallization---while being explicit that our sample cannot separate how much of the signal is selection and how much is survival. Both point the same way; disentangling them quantitatively would require a debiased, volume-limited architecture sample.

\subsection{Why GUE and not GOE? An open question}

A striking empirical fact is that the Earth systems in this program mostly sit near GOE ($\langle r\rangle\approx0.53$), whereas orbital architecture reaches GUE-plus ($0.67$). It is tempting to invoke the textbook RMT rule that GUE emerges when time-reversal symmetry is broken, and to attribute the shift to the irreversible, dissipative processes that sculpt planetary systems early on---protoplanetary gas drag, tidal damping, and orbital migration---as opposed to conservative $N$-body gravity, which is time-reversible. We raise this as a suggestive analogy and a target for theory, but we deliberately do \emph{not} claim it as established, for a reason that a careful reader will demand. The classical GOE/GUE distinction is a statement about the anti-unitary symmetry class of a \emph{Hamiltonian whose eigenvalue spectrum} is being analyzed. Orbital spacings $\Delta\ln a$ are not the eigenvalue spectrum of any such Hamiltonian, and the Earth systems we label ``GOE'' are themselves shaped by dissipation, so a literal ``dissipation $\Rightarrow$ GUE'' mapping is not self-consistent within the program. What is robust is the empirical observation: orbital repulsion is measurably more rigid (higher $\langle r\rangle$, closer to GUE) than the terrestrial charge-and-release systems. Why orbital clearing lands in a more rigid universality class---whether it reflects a genuine symmetry-class difference, the over-packing stability filter discussed above, or the specific functional form of gravitational spacing rules---is an open theoretical problem. We flag it as one of the most interesting questions the program raises, rather than resolving it by rhetorical analogy.

\subsection{The local-spacetime principle and a falsifiable prediction}

The collapse test elevates ``single source'' to ``single local spacetime frame.'' Repulsion is a property of bodies that share one gravitational well, one central mass, and one dynamical history; it does not survive the mixing of independent frames. This yields a concrete, falsifiable prediction: any sample that combines systems across genuinely different physical regimes---different stellar populations, metallicities, or formation epochs---without per-system alignment should show \emph{stronger} collapse toward Poisson than a within-population sample, because the rulers are more disparate. In the limit, were resolved orbital architectures ever measured for extragalactic systems and pooled with Galactic ones without alignment, the collapse should be sharper still. The principle is testable wherever orbital catalogs can be stratified by physical regime.

\subsection{Limitations}

\emph{Single-system sample sizes.} Target A systems are small ($n\le7$ spacings); they are qualitative, and the statistical weight rests on the pooled Target B sample. \emph{Detection biases.} The Kepler/transit sample favors compact, short-period, coplanar systems; the GUE result strictly describes the architecture of \emph{detected} multi-planet systems, which are precisely the dynamically packed ones. Whether sparser or wider systems share the same repulsion is observationally open. \emph{Unfolding choice.} Local unfolding with a smoothing spline involves a smoothing parameter; we verified that the aligned-vs-mixed contrast (0.674 vs 0.437) is robust to reasonable variation in smoothing, because it reflects frame alignment, not the detrending details. \emph{Resonant subsystems.} Beyond the flagged Galilean case, resonant chains within larger systems may locally inflate $\langle r\rangle$; the pooled result is dominated by non-resonant spacings, but a fully resonance-cleaned analysis is left for future work.

\subsection{The shape of the unified picture}

Orbital architecture extends the charge-and-release program from time into space and from the Earth to the cosmos. The mechanism is the same---a single source that depletes its local reservoir forbids closely spaced events, whether the spacing is measured in years between flares or in $\ln a$ between planets---and the diagnostic is the same $\langle r\rangle$. What the gravitational case adds is the sharpest statement yet of the program's central dichotomy: repulsion is a local-frame phenomenon, universal within a single well and erased by mixing independent ones. The gravitational shadow is the spatial face of charge-and-release, and it carves the same Wigner--Dyson rigidity into the architecture of planetary systems across the Galaxy.

\section{Conclusions}

Using textbook Solar System data and 320 Kepler/transit multi-planet systems (1130 planets) from the NASA Exoplanet Archive, we find that the logarithmic spacing of orbital semi-major axes shows GUE-class level repulsion, both for the single Solar System ($\langle r\rangle=0.632$ after Titius--Bode unfolding) and for the pooled exoplanet ensemble ($\langle r\rangle=0.674$, 95\% CI [0.68, 0.72]). The Galilean moons are correctly identified as resonance-locked (CV~$=0.09$) and excluded. Contrary to the naive expectation of Poisson collapse, pooling many systems preserves repulsion, because each system is internally repulsive and is unfolded in its own frame. The collapse is recovered only by destroying the frames: globally mixing all planets gives $\langle r\rangle=0.437$, and an independent-clock null gives 0.396. Level repulsion in orbital architecture is thus a local-spacetime property---universal within a single gravitational well, erased by the superposition of independent ones. This is the eighth system in a unified RMT program and the first to extend it from temporal to spatial, and from terrestrial to cosmic, charge-and-release.

\section*{Code and Data Availability}
All code, the reconstructed orbital dataset, and analysis outputs are at \url{https://github.com/Ruqing1963/orbital-rmt-gravitational-shadow}. Exoplanet data are from the NASA Exoplanet Archive (\texttt{pscomppars} table). The pipeline reproduces every number and figure from the raw semi-major axes.

\begin{thebibliography}{99}
\small
\bibitem[Chen(2026a)]{chen2026strat} Chen, R. (2026a). Spatial level repulsion in the stratigraphic record: a random matrix theory test of cyclostratigraphic bed spacing. Zenodo, \url{https://zenodo.org/records/20774581}.
\bibitem[Chen(2026b)]{chen2026plume} Chen, R. (2026b). Level repulsion in the deep Earth: a random matrix theory analysis of mantle-plume eruption rhythms. Zenodo, \url{https://zenodo.org/records/20768420}.
\bibitem[Chen(2026c)]{chen2026metal} Chen, R. (2026c). Charge-and-release rhythms in ore-forming systems: a random matrix theory test of metallogenic timing across four provinces. Zenodo, \url{https://zenodo.org/records/20768849}.
\bibitem[Chen(2026d)]{chen2026evo} Chen, R. (2026d). The rhythm of life and death: random matrix theory reveals level repulsion in Phanerozoic mass extinctions and radiations. Zenodo, \url{https://zenodo.org/records/20783763}.
\bibitem[Chen(2026e)]{chen2026hydro} Chen, R. (2026e). From geysers to megafloods: random matrix theory reveals universal level repulsion in hydrogeological charge-and-release systems. Zenodo, \url{https://zenodo.org/records/20780389}.
\bibitem[Chen(2026f)]{chen2026fault} Chen, R. (2026f). From hypothesis to reality: scale-dependent universality classes in Earth's seismotectonic rhythms. Zenodo, \url{https://zenodo.org/records/20768130}.
\bibitem[Chen(2026g)]{chen2026solar} Chen, R. (2026g). The depletion limit of level repulsion: why solar flares from a single active region are Poissonian. Zenodo, \url{https://zenodo.org/records/20784967}.
\bibitem[NASA Exoplanet Archive(2026)]{nasa2026} NASA Exoplanet Archive (2026). Planetary Systems Composite Parameters table. NASA/IPAC, \url{https://exoplanetarchive.ipac.caltech.edu}.
\bibitem[Bohigas et al.(1984)]{bohigas1984} Bohigas, O., Giannoni, M.-J., \& Schmit, C. (1984). Characterization of chaotic quantum spectra and universality of level fluctuation laws. \emph{Physical Review Letters}, 52, 1.
\bibitem[Atela et al.(2002)]{atela2002} Atela, P., Gol\'e, C., \& Hotton, S. (2002). A dynamical system for plant pattern formation: a rigorous analysis. \emph{Journal of Nonlinear Science}, 12, 641.
\bibitem[Hill(1878)]{hill1878} Hill, G. W. (1878). Researches in the lunar theory. \emph{American Journal of Mathematics}, 1, 5.
\end{thebibliography}

\end{document}
